\section{Empirical Validation of Concept Activation Regions}
\label{app:activation_regions}
This section provides empirical support for~\cref{def:activation} by showing that the cross-attention activations used in our formalization form coherent geometric clusters in practice.
\paragraph{Setup.}
To verify that concept activation regions are well-defined empirical objects, we extract cross-attention activations from Stable Diffusion v1.4 at the mid-block layer (\texttt{mid\_block.attentions.0}) at step $t^*{=}25$ of 50 DDIM steps for 10 concepts, each prompted with 50 shared templates differing only in the concept word.
\texttt{horse}, \texttt{pony}, \texttt{donkey}, \texttt{deer}, \texttt{dog}, \texttt{cat}, \texttt{bear}, \texttt{car}, \texttt{truck}, and \texttt{castle}.
% This visualization pool is intentionally smaller and more domain-diverse than the concept pool used for  $\hat \kappa$-estimation and CLIP classification; it is chosen to demonstrate that semantically related concepts cluster while semantically distinct concepts (e.g., vehicles vs. animals vs. architecture) separate. The $\hat \kappa$-estimator and all evaluation metrics use the full concept pool, not this visualization set.
This visualization pool is chosen to be illustrative rather than comprehensive: it deliberately spans three semantically distinct domains (animals, vehicles, architecture) so that the qualitative geometry of concept regions --- semantically related concepts clustering while distinct domains separate --- is visually apparent. Several concepts in this pool (\texttt{deer}, \texttt{car}, \texttt{truck}) are not part of the evaluation vocabulary in~\cref{tab:vocab}; they appear only to anchor the visualization's domain diversity. Conversely, $\hat\kappa$-estimation and all CLIP-based evaluation metrics use the full $65$-concept pool of~\cref{tab:vocab}, of which the visualization pool is not a strict subset.
For each concept, we generate 50 prompts using a shared set of templates that differ only in the concept word, so that differences in activation primarily reflect the concept token rather than unrelated prompt variation.

For each prompt $p$, we compute the activation
\begin{equation}
    \boldsymbol{\Phi}_{\theta}(p) \;=\; \mathrm{A}_{\theta}(\mathbf{e}_p,\ell^*,t^*) \in \mathcal{H},
\end{equation}
flatten the resulting cross-attention tensor into a vector in $\mathcal{H}$, and project all activations to two dimensions using UMAP for visualization.

\paragraph{Observations.}
\Cref{fig:umap} shows that activations induced by prompts containing the same concept form concentrated local clusters, rather than being arbitrarily scattered in activation space. This supports the view that each concept corresponds to a coherent region in the model's internal representation space.

The figure also reveals non-uniform distances between concepts. Some concepts occupy nearby regions—for example, \texttt{horse}, \texttt{pony}, and \texttt{donkey} form a tight cluster—while others, such as \texttt{castle}, are clearly separated from the animal concepts. Importantly, we do not use these visualizations to infer semantic relatedness from scratch. Rather, the figure shows that the activation geometry is structured and that concepts commonly regarded as related in the data domain tend to induce nearby activations under shared prompt templates.

These results do not by themselves prove overlap in the strict set-theoretic sense used later in Section~4. Instead, they provide empirical evidence for the weaker claim needed to motivate our definitions: concept-conditioned activations are localized, stable enough to visualize, and often organized into neighboring regions rather than isolated points. This makes it meaningful to model a concept through an activation region $\mathcal{R}_c$ and to reason about neighborhood structure in the induced activation space.

\paragraph{Choice of probe layer and timestep.}
We probe the cross-attention output at \texttt{mid\_block.attentions.0} at denoising step $t^*=25$ of $50$ throughout~\cref{app:activation_regions,app:kappa_estimation,app:displacement,app:lipschitz}, and at \texttt{up\_blocks.1.attentions.1} (same timestep) for the direct displacement measurements in~\cref{app:displacement_kappa}. The probe layer $\ell^*$ is a parameter of our framework (\cref{def:activation}), not a property of the model, and~\cref{thm:tradeoff} holds for any chosen $(\ell^*, t^*)$ at which Assumptions~\ref{asm:lip} and~\ref{asm:edit} are satisfied. We selected mid-block as the default probe based on a sensitivity analysis (\cref{app:layer_sensitivity}) over five candidate cross-attention layers spanning the U-Net depth: the $\hat\kappa$ ordering across our five target concepts is stable across all probed layers (Spearman $\rho \geq 0.80$ vs.\ mid-block), but mid-block produces the cleanest bimodal separation between high-$\hat\kappa$ and low-$\hat\kappa$ concepts and the tightest within-cluster spread among the animal subtypes. Step $t^*=25$ (the midpoint of the DDIM trajectory) corresponds to the regime in which large-scale layout and object identity are committed in latent diffusion models. The up-block layer in~\cref{app:displacement_kappa} was selected because it lies downstream of the cross-attention edits applied by fine-tuning methods and therefore reflects the cumulative effect of those edits; \cref{app:layer_sensitivity} confirms that the $\hat\kappa$ ranking at this layer is consistent with mid-block.